\subsection{Developer Onboarding}
\label{subsec:entwickler-onboarding}

Onboarding follows the sequence of cloning the master, running
\texttt{doctor}, running \texttt{init}, moving to the target project, running
\texttt{doctor} again, and then running \texttt{install} and \texttt{run}.
Product code, configuration, tests, and commits then belong in the derived
repository.
Template maintenance, by contrast, takes place in the master and validates the
entire baseline there
\parencite{kleiveist2026governance-snapshot,kleiveist2026profiles}.
This separation prevents product development from inadvertently modifying the
shared template.

The basic prerequisites are Git, Python 3.11 or later, and Node.js 24 or later.
Desktop profiles additionally require Rust Stable and platform-specific Tauri
packages; Docker, PostgreSQL, and PyGitIndex are required only for the
respective optional tasks. Optional tools therefore do not block every
profile.

\path{AGENTS.md} also provides the repository policy to supporting coding
agents. It defines 900 code LOC as a permitted absolute maximum, requires
evaluation from 600 LOC and normally a split plan above 750 LOC, prohibits
merely weakening rules, and requires thin routers, adapters, Tauri commands,
and composition roots. After structural code changes, the complete quality
gate is to run before handoff, followed by relevant tests. Changes that cross
subsystem boundaries require the complete applicable suite. Remaining findings
must be disclosed \parencite{kleiveist2026agents}.

The agent file is explicit normative context, not empirical evidence of
effectiveness. Its existence proves neither that a particular model follows
the rules reliably nor that agent-assisted code is automatically clean.
Technically enforced rules, human review, and model behavior remain separate
layers.

The README and specialized documents organize getting started, tooling,
profiles, configuration, CI, containers, and releases into separate areas.
This standardizes the public setup paths, but does not yet demonstrate a
measurable reduction in onboarding time
\parencite{kleiveist2026tooling,kleiveist2026configuration,kleiveist2026ci,kleiveist2026release}.
The documentation structure thus supplements technical onboarding, but does
not replace an empirical measurement of its duration.
\beginnerinput{workspace/statement/500_tooling_workflow/570}
